\documentclass[12pt,a4paper]{article}
\usepackage[margin=25mm, headheight=15pt, headsep=10pt]{geometry}
\usepackage{fontspec}
\usepackage[table]{xcolor} % Note: table option is required for rowcolors
\usepackage{titlesec}
\usepackage{tocloft}
\usepackage{fancyhdr}
\usepackage{booktabs}
\usepackage{tcolorbox}
\tcbuselibrary{skins, breakable}
\usepackage[colorlinks=true, linkcolor=emerald, urlcolor=emerald, bookmarks=true]{hyperref}

\definecolor{emerald}{RGB}{0,102,51}
\definecolor{darkred}{RGB}{139,0,0}
\definecolor{lightgray}{RGB}{245,245,245}

\usepackage[nil]{babel}
\babelprovide[import, main]{bengali}
\babelprovide[import]{arabic}
\babelfont{rm}[Renderer=HarfBuzz,Scale=1.0]{Noto Serif Bengali}
\babelfont{sf}[Renderer=HarfBuzz]{Noto Sans Bengali}
\babelfont{tt}[Renderer=HarfBuzz]{Noto Sans Bengali}
\babelfont[arabic]{rm}[Renderer=HarfBuzz,Scale=1.1]{Noto Naskh Arabic}

% Section styling
\titleformat{\section}{\Large\bfseries\color{emerald}}{\thesection.}{0.5em}{}
\titleformat{\subsection}{\large\bfseries\color{emerald!85!black}}{\thesubsection.}{0.5em}{}
\titleformat{\subsubsection}{\normalsize\bfseries\color{emerald!70!black}}{\thesubsubsection.}{0.5em}{}
\titlespacing*{\section}{0pt}{18pt}{8pt}

% Fancy Header/Footer setup
\pagestyle{fancy}
\fancyhf{} % clear all header and footer fields
\fancyhead[L]{\color{emerald}\small\textbf{\nouppercase{\leftmark}}}
\fancyfoot[L]{\scriptsize\color{black!50}{\lat{AI}}-সংকলিত, আলেম-যাচাইকৃত নয়}
\fancyfoot[C]{\color{emerald!70!black}\thepage}
\renewcommand{\headrulewidth}{0.4pt}
\renewcommand{\headrule}{\hbox to\headwidth{\color{emerald}\leaders\hrule height \headrulewidth\hfill}}
\renewcommand{\footrulewidth}{0pt}

% Custom boxes
% 1. Ayat/Hadith Box
\newtcolorbox{ayatbox}[1][]{
    enhanced, breakable, 
    colback=emerald!5, 
    colframe=emerald, 
    boxrule=0.5pt, 
    arc=3pt, 
    left=10pt, right=10pt, top=10pt, bottom=10pt,
    #1
}

% 2. Quote/Translation Box
\newtcolorbox{quotebox}[1][]{
    enhanced, breakable,
    frame hidden,
    colback=lightgray,
    borderline west={3pt}{0pt}{emerald},
    left=15pt, right=10pt, top=10pt, bottom=10pt,
    sharp corners,
    #1
}

% 3. AI-generated notice box
\newtcolorbox{warnbox}[1][]{
    enhanced, breakable,
    colback=red!4,
    colframe=darkred,
    boxrule=0.8pt,
    arc=3pt,
    left=10pt, right=10pt, top=10pt, bottom=10pt,
    #1
}

% Commands
\newcommand{\arab}[1]{\foreignlanguage{arabic}{#1}}
\newcommand{\kib}[1]{\textcolor{emerald}{\textbf{#1}}}

% Latin (English) fallback: Noto Serif Bengali has NO Latin glyphs
\newfontfamily\latfont[Renderer=HarfBuzz]{Noto Serif}
\newcommand{\lat}[1]{{\latfont #1}}

\setlength{\parskip}{5pt}
\setlength{\parindent}{0pt}

% Fix for missing bullet in Noto Serif Bengali
\renewcommand{\labelitemi}{$\bullet$}



\begin{document}
\begin{center}{\Large\bfseries\color{emerald} কুরআনের আয়াত — ছবি আঁকার বিধানের প্রাসঙ্গিক আয়াত}\end{center}
\vspace{1em}
\section*{কুরআনের আয়াত — ছবি আঁকার বিধানের প্রাসঙ্গিক আয়াত}

\begin{warnbox}
 \textbf{গুরুত্বপূর্ণ নোটিশ (\lat{Important} \lat{Notice}):} এই কন্টেন্টটি \lat{AI} (কৃত্রিম বুদ্ধিমত্তা) দিয়ে প্রস্তুত ও সংকলিত। \lat{AI} কঠোর সূত্র-মান নীতি মেনে শুধু নির্ভরযোগ্য উৎস থেকে তথ্য সংগ্রহ করেছে — কেবল সহীহ/হাসান হাদিস ও সহীহ সনদের আসার রাখা হয়েছে, দুর্বল সনদ বাদ দেওয়া হয়েছে। তবে এটি কোনো আলেম বা মুহাদ্দিস কর্তৃক যাচাইকৃত নয়।
\end{warnbox}

\begin{quote}
\textbf{গুরুত্বপূর্ণ সতর্কতা (\lat{Important} \lat{Note}):} কুরআনে \textbf{সরাসরি} \lat{“}ছবি আঁকা হারাম\lat{”} বলা হয়নি। ছবি আঁকার নিষেধের মূল ভিত্তি \textbf{সহীহ হাদিস}। নিচের আয়াতগুলো সেই নিষেধের অন্তর্নিহিত \textbf{\lat{reasoning} (যৌক্তিকতা)} ও প্রেক্ষাপট বোঝায় — কেন ইসলাম ছবি-মূর্তিকে এত সতর্কতার সাথে দেখে। তাই হাদিস ফাইলটি এই ফাইলের সাথে একসাথে পড়ুন। \\
তাফসীর সংগ্রহ করা হয়েছে প্রামাণ্য গ্রন্থ থেকে: তাফসীর ইবনে কাসীর, জামিউল বায়ান (তাবারী), তাফসীরুল কুরতুবী, তাফসীর আস-সা'দী, মাআরিফুল কুরআন (মুফতী মুহাম্মদ শাফী উসমানী)।
\end{quote}

\medskip\hrule\medskip

\subsection*{আয়াত ১ — সূরা আল-আরাফ (৭) — আয়াত ১৯১}

\subsubsection*{১. সূত্র কার্ড}

\begin{table}[h]\centering\small
\begin{tabular}{ll}
বিষয় & বিবরণ \\
\hline
\textbf{সূরা} & আল-আরাফ (৭) — মক্কী \\
\textbf{আয়াত} & ১৯১ \\
\textbf{পারা} & ৯ \\
\textbf{মূল বিষয়} & যে মূর্তি/ছবি নিজে সৃষ্ট নয়, সে কিছুই সৃষ্টি করতে পারে না — তাই আল্লাহর সাথে শরিক হতে পারে না \\
\textbf{প্রসঙ্গ} & মুশরিকদের মূর্তিপূজার খণ্ডন \\
\hline
\end{tabular}
\end{table}

\subsubsection*{২. প্রসঙ্গ ও প্রেক্ষাপট (\lat{Context})}

সূরা আরাফের এই অংশে আল্লাহ মুশরিকদের প্রশ্নের উত্তর দিচ্ছেন — যারা আল্লাহর সাথে এমন কিছু শরিক করে যাদের কোনো \textbf{\lat{creative} \lat{power} (সৃষ্টিশক্তি)} নেই। আয়াত ১৯০-১৯২ একই ধারাবাহিকতা। মুশরিকরা পাথর-কাঠের মূর্তি বানিয়ে তাদের আল্লাহর অংশীদার ভাবত; আল্লাহ বলছেন — যারা নিজেরাই সৃষ্ট, তারা তো কিছুই সৃষ্টি করতে পারে না। এই আয়াতই বোঝায় — ছবি ও মূর্তির আশেপাশে ইসলামের \textbf{সংবেদনশীলতা (\lat{sensitivity})} কোথা থেকে আসে: মানুষের তৈরি ছবি-মূর্তি কখনো সৃষ্টিকর্তার প্রতিদ্বন্দ্বী হতে পারে না, বরং বান্দার জন্য এটি বিভ্রান্তির কারণ।

\subsubsection*{৩. মূল আরবি}

\begin{center}\arab{أَيُشْرِكُونَ مَا لَا يَخْلُقُ شَيْئًا وَهُمْ يُخْلَقُونَ}\end{center}
\textbf{উচ্চারণ:} আ-ইউশরিকূনা মা- লা- ইয়াখলুকু শাইআন ওয়া হুম ইউখলাকূন।

\subsubsection*{৪. বাংলা অর্থ (\lat{pure})}

\textbf{\lat{“}তারা কি (আল্লাহর সাথে) এমন শরিক করে, যারা কিছুই সৃষ্টি করে না, অথচ তারা নিজেরাই সৃষ্ট?\lat{”}}

\subsubsection*{৫. শব্দের ব্যাখ্যা (\lat{Vocabulary})}

\begin{table}[h]\centering\small
\begin{tabular}{ll}
শব্দ & অর্থ \\
\hline
\textbf{আ-ইউশরিকূন (\arab{أيشركون})} & তারা কি শরিক করে? — শিরক করার প্রশ্নসূচক নিন্দা \\
\textbf{ইয়াখলুকু (\arab{يخلق})} & সৃষ্টি করা — শুধু আল্লাহর কাজ \\
\textbf{ওয়া হুম ইউখলাকূন} & অথচ তারা নিজেরাই সৃষ্ট — নিজেরা বানানো, তাই বানানো যায় না সৃষ্টি \\
\hline
\end{tabular}
\end{table}

\subsubsection*{৬. গভীর তাফসীর (\lat{Deep} \lat{Tafsir})}

\begin{itemize}
\item \textbf{ইবনে কাসীর:} এখানে আল্লাহ মুশরিকদের অপরিণামদর্শিতা দেখিয়েছেন — তারা এমন বস্তুকে অংশীদার বানিয়েছে যে নিজে সৃষ্ট হওয়া সত্ত্বেও কিছুই সৃষ্টি করতে পারে না। সৃষ্টির ক্ষমতা একমাত্র আল্লাহর। (তাফসীর ইবনে কাসীর, ৩/৪৭৪)
\item \textbf{সা'দী:} আয়াতটি শিরকের মূল দুর্বলতা তুলে ধরেছে — শরিক যদি নিজেই সৃষ্ট, তবে সে আল্লাহর মতো সৃষ্টি করবে কীভাবে? ছবি-মূর্তি বানানো নিষেধের কারণও এখান থেকে বোঝা যায়: মানুষকে \lat{“}সৃষ্টিকর্তা-অনুকরণ\lat{”} থেকে বিরত রাখা। (তাফসীর আস-সা'দী, পৃষ্ঠা ৩১৮)
\item \textbf{কুরতুবী:} মুশরিকদের মূর্তি তাদের কাছে সবচেয়ে বড় \textbf{\lat{honour} (সম্মান)} ছিল; আল্লাহ তা নস্যাত করছেন — কারণ এসব কিছুই বানাতে পারে না। (তাফসীরুল কুরতুবী, ৭/৩৩৫)
\end{itemize}
\subsubsection*{৭. সংশ্লিষ্ট হাদিস ও আয়াত}

\begin{itemize}
\item \textbf{আয়াত (হজ্জ ২২:৭৩):} \lat{“}তারা একটি মাছিও সৃষ্টি করতে পারবে না\lat{”} — একই বক্তব্যের পরের ধাপ।
\item \textbf{আয়াত (নাহল ১৬:১৭):} \lat{“}যিনি সৃষ্টি করেন, তিনি কি তাদের মতো, যারা সৃষ্টি করে না?\lat{”}
\item \textbf{হাদিস (বুখারি ৫৯৫১):} \lat{“}যারা ছবি বানায়, কিয়ামতে তাদের বলা হবে — যা সৃষ্টি করেছ তাতে প্রাণ দাও।\lat{”}
\end{itemize}
\subsubsection*{১. সূত্র কার্ড}

\begin{table}[h]\centering\small
\begin{tabular}{ll}
বিষয় & বিবরণ \\
\hline
\textbf{সূরা} & আল-হজ্জ (২২) — মাদানী \\
\textbf{আয়াত} & ৭৩ \\
\textbf{পারা} & ১৭ \\
\textbf{মূল বিষয়} & আল্লাহর বদলে যাদের ডাকা হয়, তারা একটি মাছিও সৃষ্টি করতে পারে না \\
\textbf{প্রসঙ্গ} & মূর্তি-উপাসনার পূর্ণ খণ্ডন \\
\hline
\end{tabular}
\end{table}

\subsubsection*{২. প্রসঙ্গ ও প্রেক্ষাপট (\lat{Context})}

সূরা হজ্জের ৭৩ নম্বর আয়াতে আল্লাহ একটি চিরন্তন \textbf{উপমা (\lat{parable})} দিয়েছেন — যারা আল্লাহকে বাদ দিয়ে মূর্তি-প্রতিমাকে ডাকে, তাদের পুরো অবস্থা একটি ছোট মাছির সামনেও অসহায়। এই আয়াতটি ছবি-মূর্তি নিয়ে ইসলামের দৃষ্টিভঙ্গির সবচেয়ে শক্তিশালী কুরআনি ভিত্তি — মূর্তি বা ছবি যতই সুন্দর হোক, তার সামর্থ্য মাছির চেয়েও কম।

\subsubsection*{৩. মূল আরবি}

\begin{center}\arab{يَا أَيُّهَا النَّاسُ ضُرِبَ مَثَلٌ فَاسْتَمِعُوا لَهُ ۚ إِنَّ الَّذِينَ تَدْعُونَ مِنْ دُونِ اللَّهِ لَنْ يَخْلُقُوا ذُبَابًا وَلَوِ اجْتَمَعُوا لَهُ ۖ وَإِنْ يَسْلُبْهُمُ الذُّبَابُ شَيْئًا لَا يَسْتَنْقِذُوهُ مِنْهُ ۖ ضَعُفَ الطَّالِبُ وَالْمَطْلُوبُ}\end{center}
\textbf{উচ্চারণ:} ইয়া আয়্যুহান নাসু দুরিবা মাছালুন ফাসতামিউ লাহ। ইন্নাল্লাযীনা তাদ'উনা মিন দূনিল্লাহি লাই ইয়াখলুকূ যুবাবাওঁ ওয়ালাও ইজতামাঊ লাহ। ওয়া ইয়াসলুবহুমুয যুবাবু শাইআন লা- ইয়াছতানকিজূহু মিনহু। দুঊফাত-তালিবু ওয়াল মাতলূব।

\subsubsection*{৪. বাংলা অর্থ (\lat{pure})}

\textbf{\lat{“}হে মানুষ! একটি উপমা পেশ করা হলো, তোমরা তা শোনো। নিশ্চয়ই আল্লাহকে বাদ দিয়ে তোমরা যাদের ডাকো, তারা কখনো একটি মাছিও সৃষ্টি করতে পারবে না — যদিও তারা সবাই এ কাজে একত্র হয়। আর মাছি যদি তাদের কাছ থেকে কিছু ছিনিয়ে নেয়, তারা তা ফিরিয়ে নিতেও সক্ষম হবে না। প্রার্থনাকারী দুর্বল এবং যাকে প্রার্থনা করা হয় সে-ও দুর্বল।\lat{”}}

\subsubsection*{৫. শব্দের ব্যাখ্যা (\lat{Vocabulary})}

\begin{table}[h]\centering\small
\begin{tabular}{ll}
শব্দ & অর্থ \\
\hline
\textbf{দুরিবা মাছালুন (\arab{ضرب} \arab{مثل})} & উপমা পেশ করা হলো \\
\textbf{লাই ইয়াখলুকূ যুবাবান (\arab{لن} \arab{يخلقوا} \arab{ذبابا})} & কখনো একটি মাছিও সৃষ্টি করতে পারবে না \\
\textbf{ইয়াসলুব (\arab{يسلب})} & ছিনিয়ে নেওয়া \\
\textbf{লাই ইয়াছতানকিজূ (\arab{لا} \arab{يستنقذوا})} & ফিরিয়ে নিতে সক্ষম হবে না \\
\textbf{দুঊফাত-তালিবু ওয়াল মাতলূব} & প্রার্থনাকারী ও প্রার্থিত বস্তু — দুই-ই দুর্বল \\
\hline
\end{tabular}
\end{table}

\subsubsection*{৬. গভীর তাফসীর (\lat{Deep} \lat{Tafsir})}

\begin{itemize}
\item \textbf{ইবনে কাসীর:} আল্লাহ সবচেয়ে দুর্বল সৃষ্টি — মাছি — এর উদাহরণ দিয়েছেন। সমস্ত মূর্তি-প্রতিমা মিলে একটি মাছিও বানাতে পারবে না; বরং একটি মাছি তাদের কাছ থেকে কিছু ছিনিয়ে নিলে তারা ফিরিয়ে পাবে না। (তাফসীর ইবনে কাসীর, ৫/৪৫৫)
\item \textbf{সা'দী:} আয়াতের শেষাংশ — \lat{“}দুর্বল প্রার্থনাকারী (মুশরিক) ও দুর্বল প্রার্থিত (মূর্তি)\lat{”} — অর্থাৎ উভয় পক্ষই অসহায়। এই আয়াত ইসলামে ছবি-মূর্তির প্রতি কঠোর দৃষ্টিভঙ্গির বীজ — কোনো প্রতিমা বা ছবির প্রতি ভক্তি-নির্ভরতা সম্পূর্ণ ভুল পথ। (তাফসীর আস-সা'দী, পৃষ্ঠা ৫৫০)
\item \textbf{মাআরিফুল কুরআন:} এটি ইসলামের শিরক-খণ্ডনের সবচেয়ে বিখ্যাত উপমা। এর শিক্ষা — সৃষ্টির সামর্থ্য নেই বলেই কোনো তৈরি বস্তু উপাস্য হতে পারে না; আর সেই বস্তুকে ঘিরে থাকা \textbf{\lat{veneration} (পূজা-ভক্তি)} সম্পূর্ণ নিষিদ্ধ। (মাআরিফুল কুরআন, ৬/৩৫৯)
\end{itemize}
\subsubsection*{৭. সংশ্লিষ্ট হাদিস ও আয়াত}

\begin{itemize}
\item \textbf{আয়াত (আরাফ ৭:১৯১):} \lat{“}যারা কিছুই সৃষ্টি করে না...\lat{”}
\item \textbf{হাদিস (বুখারি ৫৯৫০):} \lat{“}সবচেয়ে কঠিন শাস্তি ছবি নির্মাতাদের\lat{”} — হাদিসের \lat{“}নির্মাতা\lat{”} (\arab{مصور}) শব্দটি সৃষ্টির অনুকরণের ইঙ্গিত দেয়।
\end{itemize}
\subsubsection*{১. সূত্র কার্ড}

\begin{table}[h]\centering\small
\begin{tabular}{ll}
বিষয় & বিবরণ \\
\hline
\textbf{সূরা} & আন-নাহল (১৬) — মক্কী \\
\textbf{আয়াত} & ১৭ \\
\textbf{পারা} & ১৪ \\
\textbf{মূল বিষয়} & সৃষ্টিকর্তা আর অসৃষ্টিকর্তা কখনো সমান নয় \\
\textbf{প্রসঙ্গ} & মুশরিকদের মূর্তিপূজার যৌক্তিক খণ্ডন \\
\hline
\end{tabular}
\end{table}

\subsubsection*{২. প্রসঙ্গ ও প্রেক্ষাপট (\lat{Context})}

সূরা নাহলের এই আয়াতে আল্লাহ মুশরিকদের একটি সহজ কিন্তু গভীর প্রশ্ন করছেন — যে সৃষ্টি করে এবং যে সৃষ্টি করতে পারে না, তারা কি সমান? প্রশ্নটি \textbf{\lat{rhetorical} (অলঙ্কারিক)} — উত্তর স্পষ্ট: কখনো সমান নয়। মুশরিকরা যেমন মূর্তিকে আল্লাহর সাথে সমান ভাবত, ইসলাম ছবি-মূর্তির প্রতি সতর্কবার্তার যৌক্তিক ভিত্তিও এখানে।

\subsubsection*{৩. মূল আরবি}

\begin{center}\arab{أَفَمَنْ يَخْلُقُ كَمَنْ لَا يَخْلُقُ ۗ أَفَلَا تَذَكَّرُونَ}\end{center}
\textbf{উচ্চারণ:} আ-ফামাই ইয়াখলুকু কামাল লা- ইয়াখলুক। আ-ফালা- তাযাক্কারূন।

\subsubsection*{৪. বাংলা অর্থ (\lat{pure})}

\textbf{\lat{“}যিনি সৃষ্টি করেন, তিনি কি তাদের মতো, যারা সৃষ্টি করে না? তবে কি তোমরা উপদেশ গ্রহণ করবে না?\lat{”}}

\subsubsection*{৫. শব্দের ব্যাখ্যা (\lat{Vocabulary})}

\begin{table}[h]\centering\small
\begin{tabular}{ll}
শব্দ & অর্থ \\
\hline
\textbf{আ-ফামাই ইয়াখলুকু (\arab{أفمن} \arab{يخلق})} & যিনি সৃষ্টি করেন — তিনি কি? \\
\textbf{কামাল লা- ইয়াখলুক (\arab{كمن} \arab{لا} \arab{يخلق})} & তার মতো, যে সৃষ্টি করে না \\
\textbf{আ-ফালা- তাযাক্কারূন (\arab{أفلا} \arab{تذكرون})} & তবে কি তোমরা উপদেশ নেবে না? \\
\hline
\end{tabular}
\end{table}

\subsubsection*{৬. গভীর তাফসীর (\lat{Deep} \lat{Tafsir})}

\begin{itemize}
\item \textbf{ইবনে কাসীর:} আয়াতটি একটি বিচ্ছিন্ন প্রশ্ন — সৃষ্টিকর্তা ও অসৃষ্টিকর্তার মধ্যে ব্যবধান আকাশ-পাতাল। এখানে \lat{“}সৃষ্টি\lat{”} শুধু মূর্তি নয়, বরং রিজিক, জীবন-মৃত্যু, আসমান-জমিন — সবই আল্লাহর। (তাফসীর ইবনে কাসীর, ৪/৪৭৮)
\item \textbf{সা'দী:} যে ব্যক্তি নিজের \textbf{\lat{weakness} (দুর্বলতা)} সম্পর্কে সচেতন, সে কখনো সৃষ্টিকর্তাকে সৃষ্টির সাথে তুলনা করবে না। ছবি-মূর্তি নিয়ে যে ভ্রান্তি মানুষের ছিল, তার মূল — এই তুলনার ভুল। (তাফসীর আস-সা'দী, পৃষ্ঠা ৪৩৮)
\end{itemize}
\subsubsection*{৭. সংশ্লিষ্ট হাদিস ও আয়াত}

\begin{itemize}
\item \textbf{আয়াত (হজ্জ ২২:৭৩):} মাছির উপমা — একই বিষয়ের পরিপূরক।
\item \textbf{হাদিস (মুসলিম ২১১০খ):} \lat{“}যে ব্যক্তি ছবি বানায়, তাকে প্রাণ দেওয়ার জন্য বলা হবে, কিন্তু সে পারবে না\lat{”} — এই হাদিসটি এই আয়াতের জীবন্ত প্রয়োগ।
\end{itemize}
\subsubsection*{১. সূত্র কার্ড}

\begin{table}[h]\centering\small
\begin{tabular}{ll}
বিষয় & বিবরণ \\
\hline
\textbf{সূরা} & আল-আম্বিয়া (২১) — মক্কী \\
\textbf{আয়াত} & ৫১-৫৭ (নির্বাচিত: ৫২, ৫৭) \\
\textbf{পারা} & ১৭ \\
\textbf{মূল বিষয়} & ইবরাহীম (আঃ)-এর মূর্তির প্রতি সতর্ক দৃষ্টিভঙ্গি — ইসলামের চিরন্তন শিক্ষা \\
\textbf{প্রসঙ্গ} & তাওহীদের ইতিহাস — মূর্তিপূজার বিরুদ্ধে নবীদের সংগ্রাম \\
\hline
\end{tabular}
\end{table}

\subsubsection*{২. প্রসঙ্গ ও প্রেক্ষাপট (\lat{Context})}

ছবি-মূর্তির প্রতি ইসলামের মনোভাব বোঝার জন্য ইবরাহীম (আঃ)-এর ঘটনা সবচেয়ে বড় ঐতিহাসিক \textbf{\lat{context} (প্রেক্ষাপট)}। ইবরাহীম (আঃ) এমন এক জাতির মধ্যে বেড়ে উঠেছিলেন যারা হাতের তৈরি মূর্তিকে উপাস্য মানত। তার পিতা নিজেই মূর্তি তৈরি করতেন। এই আয়াতে ইবরাহীম (আঃ) সেই মূর্তিগুলোকে ঘিরে প্রশ্ন তোলেন এবং পরবর্তীতে নিজ জাতির অজান্তে মূর্তিগুলো ভেঙে দেন — তাওহীদের পক্ষে সবচেয়ে বড় ঘোষণা। এ থেকেই বোঝা যায় — ইসলামের ইতিহাসে ছবি-মূর্তির প্রতি সংবেদনশীলতা নতুন কিছু নয়, বরং নবীদের প্রাচীন সংগ্রাম।

\subsubsection*{৩. মূল আরবি}

\begin{center}\arab{إِذْ قَالَ لِأَبِيهِ وَقَوْمِهِ مَا هَٰذِهِ التَّمَاثِيلُ الَّتِي أَنْتُمْ لَهَا عَاكِفُونَ}\end{center}
\textbf{উচ্চারণ:} ইয কালা লিআবীহি ওয়া কাওমিহী মা- হাযিহিত তামাছীলুল্লাতী আনতুম লাহা- আ'কিফূন।

\begin{center}\arab{وَاللَّهِ لَأَكِيدَنَّ أَصْنَامَكُمْ بَعْدَ أَنْ تُوَلُّوا مُدْبِرِينَ}\end{center}
\textbf{উচ্চারণ:} ওয়াল্লাহি লাআকিদান্না আসনামাকুম বা'দা আন তুওয়াল্লূ মুদবিরীন।

\subsubsection*{৪. বাংলা অর্থ (\lat{pure})}

\textbf{\lat{“}যখন সে তার পিতা ও সম্প্রদায়কে বলল: এ মূর্তিগুলো কী, যাদের এবাদতে তোমরা লেগে আছ?\lat{”}} (২১:৫২)

\textbf{\lat{“}আল্লাহর কসম! তোমরা প্রস্থান করলে আমি তোমাদের মূর্তিগুলোর ব্যাপারে একটি কৌশল করব।\lat{”}} (২১:৫৭)

\subsubsection*{৫. শব্দের ব্যাখ্যা (\lat{Vocabulary})}

\begin{table}[h]\centering\small
\begin{tabular}{ll}
শব্দ & অর্থ \\
\hline
\textbf{তামাছীল (\arab{تماثيل})} & মূর্তি/প্রতিমা — ছবির সর্বোচ্চ রূপ \\
\textbf{আ'কিফূন (\arab{عاكفون})} & এবাদতে লেগে থাকা — ভক্তি নিবদ্ধ করা \\
\textbf{আসনাম (\arab{أصنام})} & মূর্তি — পূজার জন্য নির্মিত প্রতিমা \\
\textbf{কাইদ (\arab{كيد})} & কৌশল — এখানে মূর্তিভাঙার পরিকল্পনা \\
\hline
\end{tabular}
\end{table}

\subsubsection*{৬. গভীর তাফসীর (\lat{Deep} \lat{Tafsir})}

\begin{itemize}
\item \textbf{ইবনে কাসীর:} ইবরাহীম (আঃ) মূর্তির প্রতি ভক্তি-আসক্তি দেখে প্রশ্ন তুললেন — \lat{“}এসব কী, যাদের তোমরা এবাদতে লেগে আছ?\lat{”} অর্থাৎ এরা না খেতে পারে, না কথা বলতে পারে, না কোনো উপকার করতে পারে। (তাফসীর ইবনে কাসীর, ৫/৩৭৫)
\item \textbf{সা'দী:} ইবরাহীম (আঃ)-এর প্রশ্নের ধরনই ভক্তি-অন্ধকারাচ্ছন্ন জাতিকে জাগানোর জন্য ছিল — তিনি সরাসরি বললেন যে এ মূর্তিগুলো কিছুই নয়, এবং পরে প্রমাণ করলেন ভেঙে দিয়ে। (তাফসীর আস-সা'দী, পৃষ্ঠা ৫৩০)
\item \textbf{মাআরিফুল কুরআন:} এই ঘটনার শিক্ষা — নবী-রাসূলদের যুগ থেকে ইসলাম মানুষের তৈরি প্রতিমার প্রতি কোনো ভক্তি-নির্ভরতা গ্রহণ করেনি; বরং প্রতিমা ধ্বংস করাকে তাওহীদের ঘোষণা মানা হয়েছে। (মাআরিফুল কুরআন, ৬/৫২৭)
\end{itemize}
\subsubsection*{৭. সংশ্লিষ্ট হাদিস ও আয়াত}

\begin{itemize}
\item \textbf{আয়াত (সূরা আস-সাফফাত ৩৭:৯৩):} \lat{“}অতঃপর সে তাদের উপর ঘুষি মারল\lat{”} — মূর্তি ভাঙার ধাপ।
\item \textbf{হাদিস (মুসলিম ২১০৯):} \lat{“}মাক্কা বিজয়ের দিন নবীজি (সা.) কাবার আশেপাশের সব মূর্তি ভেঙে ফেলার নির্দেশ দেন\lat{”} — ইবরাহীমের পথেরই ধারাবাহিকতা (সহীহ বুখারি ৪২৮৮-এর আলোকে)।
\end{itemize}

\end{document}
